% GENERATED FILE. Edit canonical lessons, then run npm run generate:books.
% canonical-source-hash: fnv1a64:188c9ec49c9fab45
% canonical-lessons: ES-C454-obra, ES-C454-via, ES-C454-accidente, ES-C454-servicio, ES-R454-repaso-corte, ES-C454-sintesis-corte

\chapter{La vía cortada}
\label{ch:es-corte}
\hlchaptermodality{454}

\begin{chapteropening}
\textbf{By the end of this chapter:} \emph{I can follow a road or rail closure: read the barrier notice, understand a delay announcement, say what caused it and what is running instead.}

Read the message or notice this chapter's words exist for, then say the same thing back.
\end{chapteropening}

\section[la obra]{la obra — the works}

\label{lesson:ES-C454-obra}



\begin{quote}
Say \textbf{la oficina} and \textbf{arreglar}. There is a Latin word for \emph{work} hiding in the first, and it is about to come out on its own.
\end{quote}

\subsection*{You'll want to know: la obra}
\textbf{la obra} — the works: building work, and the site it happens on.

\begin{quote}
Calle cortada por \textbf{obras}.
\end{quote}

\emph{Están de obras} — they have the builders in. \emph{La obra empieza en marzo} — work starts in March. \emph{Un obrero} is a manual worker, and \emph{la mano de obra} is labour.

It also means a work of art or a play — \emph{una obra de teatro} — which is the same idea: a thing somebody made.

\begin{grammarlens}[title={Grammar Lens: obras in the plural is the roadworks}]
Singular and plural point at different things here:

\noindent
\begin{tabularx}{\linewidth}{@{}>{\raggedright\arraybackslash}X>{\raggedright\arraybackslash}X@{}}
\toprule
\textbf{Spanish} & \textbf{what it names} \\
\midrule
\emph{la obra} & the project, or the artwork \\
\emph{las obras} & the works going on, the disruption \\
\bottomrule
\end{tabularx}

Every barrier and diversion sign in Spain says \textbf{OBRAS}, plural. A Spaniard saying \emph{hay obras en mi calle} is complaining about noise and dust, not about a sculpture.

Keep \emph{estar de obras} as a phrase: it describes a building in the state of being worked on, and \emph{estar} is right because it is a condition that will pass.
\end{grammarlens}

\begin{cousinweb}
Latin \textbf{opera}, works — the plural of \textbf{opus}, a work, taken by Spanish as a singular feminine noun.

The same \emph{opus} is inside \textbf{la oficina}, which is \emph{opus} plus \emph{facere}: the place where work gets done. So a builder's site and an office are the same Latin word, arriving by different routes and landing far apart.

\emph{Opera} the music is the identical word again, borrowed through Italian: \emph{the works}. And English \textbf{operate}, \textbf{operation} and \textbf{opus} are all here too. One Latin noun for work, and it has ended up meaning a hole in the road, a room with desks, and a night at the theatre.
\end{cousinweb}

\subsection*{Your turn}
Say these aloud:

\begin{itemize}
\raggedright
  \item street closed for works
  \item they have the builders in
  \item what \emph{la obra} and \emph{la oficina} share
\end{itemize}

\begin{tcolorbox}[breakable,colback=teal!4,colframe=teal!35!black,title={Before you move on}]
The roadworks? (\textbf{Las obras}, plural.) A play? (\textbf{Una obra de teatro}.) What did Latin \emph{opus} mean? (\textbf{A work}.)
\end{tcolorbox}
\section[la vía]{la vía — the track}

\label{lesson:ES-C454-via}



\begin{quote}
Say \textbf{el andén} and \textbf{todavía}. One of those two has a road inside it, and it is not the one at the station.
\end{quote}

\subsection*{You'll want to know: la vía}
\textbf{la vía} — the track, the way, the route.

\begin{quote}
\textbf{Vía} 4. Tren con destino Valencia.
\end{quote}

\emph{No cruce la vía} — do not cross the track. \emph{La vía pública} — the public highway, the wording on every municipal notice. \emph{Por vía aérea} — by air.

The \textbf{andén} is where you stand; the \textbf{vía} is the rails in front of it. A station announcement gives you both, and mixing them up will put you on the platform edge.

\begin{grammarlens}[title={Grammar Lens: the accent that breaks a diphthong}]
\textbf{vía} is two syllables: \textbf{VÍ-a}. Without the accent, \emph{via} would be one syllable, because \textbf{i} and \textbf{a} together normally fuse.

\noindent
\begin{tabularx}{\linewidth}{@{}>{\raggedright\arraybackslash}X>{\raggedright\arraybackslash}X@{}}
\toprule
\textbf{written} & \textbf{syllables} \\
\midrule
v\textbf{ía} & VÍ-a \\
d\textbf{ía} & DÍ-a \\
fam\textbf{ilia} & fa-MI-lia \\
\bottomrule
\end{tabularx}

The mark on the \textbf{i} is not about stress position at all here. It is there to say: \textbf{this vowel refuses to join the next one}. Spanish has no other way to write that, so the accent does double duty across the language — sometimes marking stress, sometimes prising two vowels apart.

\emph{Todavía} has it for the same reason, and you have been reading it correctly for chapters without being told why.
\end{grammarlens}

\begin{cousinweb}
Latin \textbf{via}, road — the word on every Roman milestone.

And here is the payoff. \textbf{Todavía} is \emph{toda vía}: \textbf{all the way}. It meant \textquotedblleft{}by every route, in any case\textquotedblright{}, and from there slid to \textquotedblleft{}still\textquotedblright{} — the sense of something that has gone the whole distance and is \emph{still} going.

The family is everywhere once you see it: \emph{el viaje}, a journey; \emph{el viaducto}, a way carried over; English \textbf{via}, \textbf{obvious} (\emph{ob-viam}, lying in the way), \textbf{trivial} (\emph{tri-vium}, a place where three roads meet, so common it is worth nothing), and \textbf{voyage}.
\end{cousinweb}

\subsection*{Your turn}
Say these aloud:

\begin{itemize}
\raggedright
  \item track 4, and do not cross the track
  \item which one is the platform and which the rails
  \item what \emph{todavía} is made of
\end{itemize}

\begin{tcolorbox}[breakable,colback=teal!4,colframe=teal!35!black,title={Before you move on}]
The track? (\textbf{La vía}.) How many syllables? (\textbf{Two} — the accent keeps them apart.) What does \emph{todavía} literally say? (\textbf{All the way}.)
\end{tcolorbox}
\section[el accidente]{el accidente — what falls on you}

\label{lesson:ES-C454-accidente}



\begin{quote}
Say \textbf{la ocasión} and \textbf{roto}. A chance and a breakage are about to meet in one Latin verb.
\end{quote}

\subsection*{You'll want to know: el accidente}
\textbf{el accidente} — the accident.

\emph{Ha habido un accidente en la vía} — there has been an accident on the track. \emph{Un accidente de tráfico} — a road accident. \emph{Un accidente laboral} — an accident at work, the phrase on every insurance form.

Spanish says \emph{tener un accidente}, and the adverb is \textbf{por accidente} — by accident, meaning unintentionally.

\begin{grammarlens}[title={Grammar Lens: -cc- is two sounds, not one}]
\textbf{accidente} has a double \textbf{c}, and unlike English the two letters do different jobs:

\noindent
\begin{tabularx}{\linewidth}{@{}>{\raggedright\arraybackslash}X>{\raggedright\arraybackslash}X@{}}
\toprule
\textbf{letters} & \textbf{sound} \\
\midrule
first \textbf{c} (before \emph{c}) & \textbf{k} \\
second \textbf{c} (before \emph{i}) & \textbf{th} in Spain, \textbf{s} elsewhere \\
\bottomrule
\end{tabularx}

So the word is \textbf{ak-thi-DEN-te}, or \textbf{ak-si-DEN-te}. The rule underneath is one you already use: a Spanish \textbf{c} is hard before \emph{a}, \emph{o}, \emph{u} and another consonant, and soft before \emph{e} and \emph{i}.

The same double turns up in \emph{acción}, \emph{dirección} and \emph{lección}. Once you hear the \textbf{k} in front of the soft one, these stop being awkward.
\end{grammarlens}

\begin{cousinweb}
Latin \textbf{accidere}: \textbf{ad-}, towards, plus \textbf{cadere}, \textbf{to fall}. A thing that falls upon you.

And \textbf{la ocasión} is the same \emph{cadere}, with a different prefix — something that falls \textbf{to} you. Put the pair side by side and Spanish has a small map of luck:

\noindent
\begin{tabularx}{\linewidth}{@{}>{\raggedright\arraybackslash}X>{\raggedright\arraybackslash}X@{}}
\toprule
\textbf{word} & \textbf{what falls, and how} \\
\midrule
\textbf{la ocasión} & a chance falls your way \\
\textbf{el accidente} & a misfortune falls upon you \\
\bottomrule
\end{tabularx}

Neither was chosen by anybody. That is the whole idea the root carries, and it is why \emph{accidental} in English still means unplanned rather than harmful.

The same \emph{cadere} is in \textbf{caer}, to fall, and in English \textbf{cadence}, \textbf{decay} and \textbf{cascade}.
\end{cousinweb}

\subsection*{Your turn}
Say these aloud:

\begin{itemize}
\raggedright
  \item there has been an accident on the track
  \item I did it by accident
  \item how \emph{ocasión} and \emph{accidente} are related
\end{itemize}

\begin{tcolorbox}[breakable,colback=teal!4,colframe=teal!35!black,title={Before you move on}]
The accident? (\textbf{El accidente}.) What does \emph{cadere} mean? (\textbf{To fall}.) Which one lands \emph{on} you, and which drops \emph{in front of} you? (\textbf{Accidente} on you, \textbf{ocasión} in front of you.)
\end{tcolorbox}
\section[el servicio]{el servicio — the service}

\label{lesson:ES-C454-servicio}



\begin{quote}
Say \textbf{servir} and \textbf{arreglar}. Somebody comes out to the accident to do the second, and what they belong to is coming next.
\end{quote}

\subsection*{You'll want to know: el servicio}
\textbf{el servicio} — the service.

\emph{El servicio de grúa} — the recovery service, the tow truck. \emph{Los servicios de emergencia} — the emergency services. \emph{Fuera de servicio} — out of service, the sign on a broken lift or a stopped bus.

And the one that catches people out: \textbf{el servicio} on a door means \textbf{the toilet}. \emph{¿Dónde está el servicio?} is how you ask, and \emph{los servicios} is what the sign in a bar says.

\begin{grammarlens}[title={Grammar Lens: fuera de servicio, and its cousins}]
\emph{Fuera de} plus a bare noun is a fixed frame, and it turns up on notices constantly:

\noindent
\begin{tabularx}{\linewidth}{@{}>{\raggedright\arraybackslash}X>{\raggedright\arraybackslash}X@{}}
\toprule
\textbf{sign} & \textbf{what it means} \\
\midrule
\emph{fuera de servicio} & out of order \\
\emph{fuera de plazo} & past the deadline \\
\emph{fuera de temporada} & out of season \\
\bottomrule
\end{tabularx}

No article after \emph{de} in any of them. That bare noun is the signal that you are looking at a set phrase rather than a description, and trying to insert \emph{el} will sound wrong to a Spanish ear for the same reason \textquotedblleft{}out of the order\textquotedblright{} does to an English one.

Set beside \textbf{estropeado}, the difference is who is speaking: \emph{fuera de servicio} is the institution's wording, \emph{estropeado} is what a neighbour says.
\end{grammarlens}

\begin{cousinweb}
Latin \textbf{servitium}, from \textbf{servire}, to serve — and behind that \textbf{servus}, a slave.

That is the uncomfortable root of a very ordinary word, and it has not disappeared: \emph{el siervo} is a serf, and English \textbf{servile} keeps the sting. Spanish softened the rest of the family — \emph{servir} is what a waiter does, \emph{el servicio} is what a company sells — while the original sense stayed in the words nobody uses lightly.

English got \textbf{serve}, \textbf{servant} and \textbf{sergeant} from the same place. The last one is \emph{serviens}, \textquotedblleft{}one serving\textquotedblright{}, which rose from servant to soldier to rank.
\end{cousinweb}

\subsection*{Your turn}
Say these aloud:

\begin{itemize}
\raggedright
  \item the emergency services are on the way
  \item the lift is out of service
  \item how you would ask for the toilet
\end{itemize}

\begin{tcolorbox}[breakable,colback=teal!4,colframe=teal!35!black,title={Before you move on}]
Out of service? (\textbf{Fuera de servicio}.) The toilet? (\textbf{El servicio}.) What was a \emph{servus}? (\textbf{A slave}.)
\end{tcolorbox}
\section[(la obra, la vía, el accidente, el …]{(la obra, la vía, el accidente, el servicio) — from cold}

\label{lesson:ES-R454-repaso-corte}



\begin{quote}
Close the lessons before this one. Say all four: \emph{the works}, \emph{the track}, \emph{the accident}, \emph{the service}.
\end{quote}

\begin{grammarlens}[title={Grammar Lens: the plural is not always more of the same}]
Two of these four shift meaning when you add the \textbf{s}:

\noindent
\begin{tabularx}{\linewidth}{@{}>{\raggedright\arraybackslash}X>{\raggedright\arraybackslash}X@{}}
\toprule
\textbf{singular} & \textbf{plural} \\
\midrule
\emph{la obra} — the project & \emph{las obras} — the disruption \\
\emph{el servicio} — the service & \emph{los servicios} — the toilets \\
\bottomrule
\end{tabularx}

This is worth holding because the plural is the form you will actually read. Every diversion sign says \textbf{OBRAS} and every bar door says \textbf{SERVICIOS}, so the singular meanings you learned are the ones you will meet less often.

\emph{La vía} and \emph{el accidente} behave themselves and mean exactly more of the same thing.
\end{grammarlens}

\begin{grammarlens}[title={Grammar Lens: every one of the four was already in the corpus}]
\noindent
\begin{tabularx}{\linewidth}{@{}>{\raggedright\arraybackslash}X>{\raggedright\arraybackslash}X@{}}
\toprule
\textbf{new word} & \textbf{where its root already was} \\
\midrule
\textbf{la obra} & \emph{la oficina} — \emph{opus} + \emph{facere} \\
\textbf{la vía} & \emph{todavía} — \emph{toda vía}, all the way \\
\textbf{el accidente} & \emph{la ocasión} — both from \emph{cadere}, to fall \\
\textbf{el servicio} & \emph{servir} \\
\bottomrule
\end{tabularx}

\emph{Todavía} is the sharpest of the four. It is among the commonest words in the language, you have been using it for chapters, and the road was sitting inside it the whole time.

That keeps happening — \emph{mira} gave \emph{mirar}, \emph{sí claro} will give \emph{claro}, \emph{el punto} gave \emph{apuntar}. When a new word looks like a piece of something you already say, check before learning it twice.
\end{grammarlens}

\subsection*{Your turn}
Say these aloud:

\begin{itemize}
\raggedright
  \item the street is closed for works
  \item do not cross the track
  \item which two words change meaning in the plural
\end{itemize}

\begin{tcolorbox}[breakable,colback=teal!4,colframe=teal!35!black,title={Before you move on}]
All four again, in Spanish, without looking. Then one sentence that closes a road.
\end{tcolorbox}
\section[Practice]{(calle cortada) — read it, then do it yourself}

\label{lesson:ES-C454-sintesis-corte}



\begin{quote}
Say \textbf{las obras} and \textbf{la vía}. Each one closes something below.
\end{quote}

\subsection*{Your turn}
Two notices on the same morning. Read both, then answer.

\textbf{1 — on a barrier}

\begin{quote}
Calle cortada por \textbf{obras} hasta el día 12. Paso de peatones por la acera de enfrente.
\end{quote}

\textbf{2 — over the loudspeaker}

\begin{quote}
Atención. Por un \textbf{accidente} en la \textbf{vía}, el tren con destino Valencia saldrá con retraso. El \textbf{servicio} de autobús sustituirá al tren hasta las once.
\end{quote}

Say these aloud:

\begin{itemize}
\raggedright
  \item how long the street is shut, and where to walk
  \item why the train is late
  \item what is replacing it, and until when
\end{itemize}

\begin{culture}
\textbf{Por un accidente en la vía} is doing the work of an English clause — \textquotedblleft{}because there has been an accident on the line\textquotedblright{}. \emph{Por} plus a noun is how Spanish gives a cause on a notice, and it is shorter than any sentence you could build with \emph{porque}. The same shape opened the first notice: \emph{cortada por obras}.

Then notice \textbf{el servicio de autobús sustituirá al tren}. The announcement names the \emph{service}, not the buses, and the service is what acts. An English announcement would more likely say \textquotedblleft{}buses will replace the train\textquotedblright{}. Spanish prefers the institution as the subject in this register, which is the opposite of what the delivery card did — and the difference is that a card is apologising while an announcement is informing.

\textbf{Cortada} agrees with \emph{la calle}: feminine, because the street is.
\end{culture}

\begin{tcolorbox}[breakable,colback=teal!4,colframe=teal!35!black,title={Before you move on}]
Now your turn, out loud and then in writing. Four lines: say the street is closed for works, say until when, say there has been an accident on the track, and say a bus service is replacing the train.

Then check yourself: did you use \emph{por} plus a noun to give the reason, and did \emph{cortada} agree with what it describes?
\end{tcolorbox}
